
\subsection{Module 4 — Fusions (Mergers)}

Ce module implémente la détection automatique de cannibalisation inter-agences. L'algorithme \texttt{cKDTree} identifie toutes les paires d'agences situées à moins de 1.2 km l'une de l'autre. Pour chaque paire détectée, le module affiche :

\begin{itemize}
    \item Les noms et scores composites des deux agences.
    \item La distance exacte séparant les deux points (en km).
    \item L'action recommandée pour chacune : conserver la plus performante, fermer ou fusionner la moins rentable.
\end{itemize}

Le résultat est présenté sous forme de tableau interactif triable et d'une carte Folium dédiée affichant les paires reliées par des lignes rouges, permettant une visualisation géographique immédiate du phénomène de cannibalisation.

\begin{figure}[h]
    \centering
    \fbox{\parbox{0.9\textwidth}{\centering\vspace{3cm}\textit{Capture d'écran : Détection des paires cannibales avec distances et recommandations.}\vspace{3cm}}}
    \caption{Module Mergers : Identification automatique des agences trop proches avec recommandation de fusion.}
    \label{fig:screenshot_mergers}
\end{figure}

\subsection{Module 5 — Analyse par Ville (Cities)}

Le module Cities fournit une analyse macroscopique agrégée par agglomération. Pour chaque ville du réseau, le système calcule et affiche :

\begin{itemize}
    \item Le nombre total d'agences dans la ville.
    \item Le score composite moyen, minimum et maximum.
    \item Le score commercial moyen et le score de compétition moyen.
    \item La proportion d'agences classées KEEP vs CLOSE.
\end{itemize}

Un graphique à barres Plotly permet de comparer visuellement les villes entre elles. L'utilisateur peut sélectionner une ville spécifique pour obtenir un zoom détaillé sur ses agences individuelles.

\begin{figure}[h]
    \centering
    \fbox{\parbox{0.9\textwidth}{\centering\vspace{3cm}\textit{Capture d'écran : Analyse comparative des villes avec métriques agrégées.}\vspace{3cm}}}
    \caption{Module Cities : Tableau de bord macroscopique par agglomération.}
    \label{fig:screenshot_cities}
\end{figure}

\subsection{Module 6 — Audit Individuel (Deep Dive)}

Ce module constitue la fiche d'identité 360° d'une agence. L'utilisateur sélectionne une agence via un menu déroulant (\texttt{st.selectbox}) et obtient instantanément :

\begin{enumerate}
    \item \textbf{Jauge Composite Score :} Un indicateur Plotly \texttt{Gauge} affichant le score composite avec des zones colorées (rouge $< 0.20$, jaune $0.20-0.35$, vert $> 0.35$) et un delta par rapport à la moyenne du réseau.
    \item \textbf{Scores Détaillés :} Les 6 scores spatiaux individuels présentés sous forme de métriques \texttt{st.metric()} avec indicateurs de tendance.
    \item \textbf{Informations Contextuelles :} Le cluster d'appartenance, l'action recommandée, les coordonnées GPS, et la ville.
    \item \textbf{Carte de Localisation :} Une mini-carte Folium centrée sur l'agence sélectionnée avec un rayon de 1km illustrant la zone d'influence.
\end{enumerate}

\begin{figure}[h]
    \centering
    \fbox{\parbox{0.9\textwidth}{\centering\vspace{3cm}\textit{Capture d'écran : Fiche 360° d'une agence avec jauge, scores et carte.}\vspace{3cm}}}
    \caption{Module Deep Dive : Audit individuel complet d'une agence sélectionnée.}
    \label{fig:screenshot_deepdive}
\end{figure}

\subsection{Module 7 — Données Brutes (Data)}

Ce module offre une vue transparente sur les données sous-jacentes. Il permet aux analystes de vérifier, filtrer et exporter les datasets utilisés par les algorithmes. Le tableau interactif \texttt{st.dataframe()} supporte le tri par colonne, le filtrage textuel, et la sélection de colonnes à afficher. Les colonnes numériques (comme \texttt{composite\_score}) sont enrichies de barres de progression (\texttt{st.column\_config.ProgressColumn}) pour une lecture visuelle immédiate.

\subsection{Module 8 — Site Selection (Benchmarking Statique)}

Ce module implémente une première approche de sélection de site par benchmarking KNN statique. Pour un ensemble de candidats pré-chargés, l'algorithme calcule la similarité cosinus entre le vecteur spatial de chaque candidat et les centroïdes des 5 clusters K-Means. Les 5 agences existantes les plus similaires ``votent'' sur le potentiel du candidat, fournissant une estimation de revenu attendu.

\begin{figure}[h]
    \centering
    \fbox{\parbox{0.9\textwidth}{\centering\vspace{3cm}\textit{Capture d'écran : Module Site Selection avec benchmarking des candidats.}\vspace{3cm}}}
    \caption{Module Site Selection : Évaluation statique des sites candidats par similarité cosinus.}
    \label{fig:screenshot_sitesel}
\end{figure}

\subsection{Module 9 — Run Site Selection (Moteur Interactif)}

Ce module constitue le moteur de recommandation le plus avancé de l'application. Contrairement au module 8 (statique), il permet à l'utilisateur de paramétrer dynamiquement l'ensemble de la recherche via une série de contrôles interactifs :

\begin{enumerate}
    \item \textbf{Upload de Données :} L'utilisateur charge son propre fichier CSV de branches scorées et un fichier de données POI fusionnées.
    \item \textbf{Paramètres Spatiaux (Sliders) :}
    \begin{itemize}
        \item Distance minimum aux agences existantes (km) : contrôle anti-cannibalisation.
        \item Rayon d'agrégation POI (km) : taille de la zone d'influence.
        \item Nombre de clusters K-Means : granularité de la segmentation.
        \item Top N candidats à retourner.
    \end{itemize}
    \item \textbf{Pipeline Automatisé :} Au clic du bouton ``Run'', le système exécute séquentiellement :
    \begin{enumerate}
        \item Chargement et validation des données.
        \item Construction du cKDTree pour l'enrichissement en densité de population.
        \item Clustering K-Means et identification du profil ``High Performer''.
        \item Génération de la grille de candidats (un point tous les 500m).
        \item Calcul des scores POI par agrégation radiale via cKDTree.
        \item Filtrage par contraintes (distance minimum, densité, diversité par ville).
        \item Calcul de la similarité cosinus et sélection des Top N.
    \end{enumerate}
    \item \textbf{Visualisation des Résultats :} Les candidats sélectionnés sont affichés sur une carte interactive ScatterMapbox avec tooltips détaillés, accompagnés d'un tableau récapitulatif exportable.
\end{enumerate}

\begin{figure}[h]
    \centering
    \fbox{\parbox{0.9\textwidth}{\centering\vspace{3cm}\textit{Capture d'écran : Moteur de Site Selection avec sliders, carte des candidats et tableau des résultats.}\vspace{3cm}}}
    \caption{Module Run Site Selection : Moteur interactif de découverte de nouveaux emplacements avec contraintes paramétrables.}
    \label{fig:screenshot_runsite}
\end{figure}

\subsection{Module 10 — Optimisation Réseau (Network Optimization)}

Ce module implémente la Matrice BCG Spatiale décrite en détail au Chapitre 4. L'interface se compose de trois sections principales :

\begin{enumerate}
    \item \textbf{Configuration :} L'utilisateur paramètre les seuils de sous-performance et de sur-performance via des sliders. Ces seuils contrôlent la frontière entre les quadrants FLAGSHIP, REVIEW/INVEST, MONITOR et CLOSE.
    \item \textbf{Matrice BCG Interactive :} Un scatter plot Plotly affichant chaque agence dans l'espace (Potentiel Prédit vs Revenu Réel). Les quadrants sont colorés, et le survol d'un point affiche les détails complets de l'agence. Les lignes médianes et la diagonale d'espérance ($y = x$) sont tracées en superposition.
    \item \textbf{Tableau d'Actions :} Un résumé tabulaire des recommandations stratégiques, filtré par quadrant. Les agences REVIEW/INVEST et CLOSE sont mises en évidence avec le détail de leur écart (gap) entre potentiel et performance.
\end{enumerate}

Le module intègre également la détection automatique des agences FLAGSHIP proches d'agences à fermer, suggérant des opportunités de consolidation d'actifs.

\begin{figure}[h]
    \centering
    \fbox{\parbox{0.9\textwidth}{\centering\vspace{3cm}\textit{Capture d'écran : Matrice BCG interactive avec quadrants colorés et tableau d'actions.}\vspace{3cm}}}
    \caption{Module Network Optimization : Matrice BCG spatiale interactive avec recommandations stratégiques automatisées.}
    \label{fig:screenshot_bcg}
\end{figure}

\subsection{Module 11 — Enrichissement de Données (Data Enrichment)}

Ce dernier module transforme l'application d'un outil d'analyse statique en une plateforme auto-alimentée. Il permet à l'utilisateur de lancer directement le scraper Google Maps depuis l'interface du dashboard.

\textbf{Architecture Technique :} L'onglet implémente un paradigme producteur-consommateur asynchrone :
\begin{itemize}
    \item Le scraper est lancé comme processus enfant via \texttt{subprocess.Popen} avec \texttt{PYTHONIOENCODING=utf-8}.
    \item Un thread daemon lit le \texttt{stdout} en continu et alimente une \texttt{queue.Queue} thread-safe.
    \item La boucle Streamlit vide la queue et affiche la sortie dans un bloc \texttt{st.code()} simulant un terminal en temps réel.
\end{itemize}

\textbf{Contrôles Utilisateur :}
\begin{itemize}
    \item \textbf{Mode de Scraping :} Choix entre ``Reprendre depuis les checkpoints'' ou ``Recommencer à zéro''.
    \item \textbf{Boutons Start/Stop :} Lancement et arrêt gracieux du scraper (les données sont sauvegardées).
    \item \textbf{Console Live :} Affichage en temps réel de la progression exacte du scraper (nombre de POI trouvés, cellule en cours, etc.), reproduisant fidèlement la sortie du terminal.
    \item \textbf{Pipeline Clean \& Merge :} Un bouton dédié exécute séquentiellement \texttt{clean.py} puis \texttt{cleanall.py}, intégrant les nouvelles données dans \texttt{merged\_data\_clean.csv}.
\end{itemize}

\begin{figure}[h]
    \centering
    \fbox{\parbox{0.9\textwidth}{\centering\vspace{3cm}\textit{Capture d'écran : Console de scraping en direct avec contrôles et pipeline de nettoyage.}\vspace{3cm}}}
    \caption{Module Data Enrichment : Interface de scraping live avec console temps réel et pipeline automatisé de nettoyage et fusion.}
    \label{fig:screenshot_enrichment}
\end{figure}

\section{Valeur Stratégique et Apport de la Solution}

La plateforme Branch Intelligence apporte une valeur stratégique mesurable à Attijariwafa Bank sur quatre axes :

\subsection{Rationalisation des Décisions d'Investissement}

Avant ce projet, les décisions d'ouverture et de fermeture d'agences reposaient sur l'intuition des directeurs régionaux et des études de marché ponctuelles. La matrice BCG spatiale (Module 10) fournit désormais un cadre quantitatif objectif : les 86 agences identifiées comme REVIEW/INVEST représentent un potentiel de revenus inexploité, tandis que les 513 agences CLOSE constituent des cibles de rationalisation budgétaire.

\subsection{Réduction de la Cannibalisation}

Le module Mergers (Tab 4) a identifié automatiquement les paires d'agences se disputant le même bassin de clientèle. La consolidation de ces doublons permettrait de réduire les coûts opérationnels tout en maintenant la couverture géographique du réseau.

\subsection{Expansion Scientifique du Réseau}

Le moteur de Site Selection (Module 9) remplace les études de faisabilité traditionnelles par un algorithme reproductible et paramétrable. L'utilisation de la similarité cosinus garantit que les nouveaux sites recommandés reproduisent le profil spatial des agences les plus performantes du réseau, maximisant la probabilité de succès commercial.

\subsection{Autonomie Opérationnelle}

Le module Data Enrichment (Tab 11) assure la pérennité du système. L'équipe Analytics peut enrichir la base de données POI de manière autonome, sans intervention de l'équipe de développement, puis relancer les analyses avec des données fraîches. Cette boucle d'auto-alimentation transforme le projet d'un livrable ponctuel en une plateforme vivante.

\section{Conclusion}

La réalisation de la plateforme Branch Intelligence démontre la faisabilité technique et la valeur métier d'une approche Data-Driven appliquée à l'optimisation du réseau bancaire. L'application, composée de plus de 3\,000 lignes de code Python et structurée en 11 modules fonctionnels, couvre l'intégralité du cycle de vie de la donnée géospatiale : acquisition, nettoyage, modélisation, visualisation et aide à la décision.

Les résultats empiriques présentés au Chapitre 4 --- un $R^2$ de 0.88 en régression et une précision de 82.9\% en classification --- valident la robustesse des modèles intégrés. La matrice BCG spatiale, innovation centrale du projet, permet pour la première fois de quantifier objectivement l'écart entre le potentiel géographique d'un emplacement et la performance managériale de l'équipe en place.

L'architecture asynchrone du module de scraping et le système de checkpoints garantissent la résilience opérationnelle du système, tandis que le choix de Streamlit assure un déploiement rapide et une maintenance simplifiée. La plateforme est prête pour un déploiement pilote auprès des équipes de planification stratégique d'Attijariwafa Bank.
